% Generated draft: requires mathematical and source review.
\documentclass[11pt]{article}
\usepackage[margin=25mm]{geometry}
\usepackage{fontspec}
\setmainfont{Latin Modern Roman}
\usepackage{amsmath,amssymb,mathtools,unicode-math}
\setmathfont{Latin Modern Math}
\usepackage{enumitem,needspace,xurl,hyperref}
\setlength{\parindent}{0pt}
\setlength{\parskip}{5pt}
\setlength{\emergencystretch}{3em}
\newcommand{\N}{\mathbb{N}}
\newcommand{\Z}{\mathbb{Z}}
\newcommand{\Q}{\mathbb{Q}}
\newcommand{\R}{\mathbb{R}}
\newcommand{\C}{\mathbb{C}}
\newcommand{\F}{\mathbb{F}}
\newcommand{\A}{\mathbb{A}}
\newcommand{\PP}{\mathbb{P}}
\newcommand{\E}{\mathbb{E}}
\newcommand{\eps}{\varepsilon}
\newcommand{\dd}{\,\mathrm{d}}
\begin{document}
\section*{UnsolvedMath ID: 30002370}
\textbf{Source problem code:} catalogue entry
\subsection{Definitions and mathematical statement}
An Einstein metric on a smooth manifold \(M\) is a Riemannian metric \(g\) satisfying
\[
\operatorname{Ric}(g)=\lambda g
\]
for a constant \(\lambda\in\R\). The supplied record does not specify the dimension, the Einstein constants, or whether the pieces are compact, complete, or subject to boundary conditions. For compatible smooth manifolds of the same dimension, a connected sum \(M_1\#M_2\) is formed by removing the interiors of embedded balls and identifying the resulting boundary spheres with suitable gluing data. The record does not specify whether its limit space or bubble is itself such a smooth connected-sum piece.

For a differentiable functional \(\mathcal F\) on a specified class \(\mathcal A(M)\) of Riemannian metrics, a critical metric \(g\in\mathcal A(M)\) satisfies
\[
\left.\frac{\dd}{\dd t}\right|_{t=0}\mathcal F(g_t)=0
\]
for every admissible smooth variation \(g_t\in\mathcal A(M)\) with \(g_0=g\). Both the functional and the admissible class, including any normalization or boundary conditions, must be specified to determine this equation.

\textbf{Source question, with unresolved specifications.} Can a reverse-bubbling construction, starting with a limit space and attaching a bubble, glue Einstein metrics on connected-sum pieces to produce a critical metric of a Riemannian functional generalizing the Einstein--Hilbert functional? In a smooth two-piece setting, the desired conclusion would have the schematic form
\[
 g\in\mathcal A(M_1\#M_2),\qquad
 \left.\frac{\dd}{\dd t}\right|_{t=0}\mathcal F(g_t)=0
 \quad\text{for every admissible variation through }g,
\]
with \(g\) obtained by the proposed gluing construction. This display specifies the meaning of the requested criticality; it does not assert existence or select the missing gluing hypotheses.

The source does not identify \(\mathcal F\), define the permitted limit spaces or bubbles, prescribe the attachment and convergence conditions, or state whether existence is sought for every admissible input or only for particular examples. Consequently, a fully quantified existence assertion cannot be extracted from this record without additional source information. In particular, the requested output is a critical metric of the unspecified generalized functional; the record does not require that output metric itself to be Einstein.
\subsection{Short English statement}
The question asks whether attaching a bubble to a limit space can glue Einstein pieces into a connected-sum metric that is critical for a functional generalizing the Einstein--Hilbert functional. The functional, admissible geometric data, and scope of the existence claim remain unspecified, so this is a qualified statement-only draft requiring mathematical review.
\subsection{Sources}
\begin{itemize}
\item Local UnsolvedMath registry, \texttt{data/unsolvedmath/problems.json}, permanent ID \(30002370\), source problem code \texttt{catalogue entry}, title \emph{Einstein Metric Gluing by Reverse Bubbling}; supplied complete source record from the Oberwolfach Reports Open Problems collection. No individual proposer is identified.
\item Supplied report citation: \emph{Differentialgeometrie im Großen (2014)}, Oberwolfach Reports, DOI \texttt{10.4171/owr/2013/33}, pp.\ 1929--1974. \url{https://doi.org/10.4171/owr/2013/33}. The registry records proposed year 2013; the relationship between that year, the DOI year component, and the title's 2014 date requires source verification. No problem-specific page or named contributor is supplied.
\item Supplied literature reference: Gursky and Viaclovsky, ``gluing critical metrics.'' \url{https://arxiv.org/abs/1303.0827}. This is the supplied descriptive label; a full bibliographic citation is not provided.
\item Supplied literature reference: Gursky and Viaclovsky, ``connected-sum metrics.'' \url{https://doi.org/10.1016/j.aim.2015.11.054}. This is the supplied descriptive label; a full bibliographic citation is not provided.
\item Literature assessment dated exactly \texttt{2026-08-21}: the supplied curation labels the problem partially solved and reports that Gursky and Viaclovsky construct critical metrics on several connected sums by gluing, while the general reverse-bubbling existence question remains unsolved. This is supplied dated curation derived from the OWR collection, not an independently verified assessment; its time-sensitive status and cited claims require independent verification.
\item The registry labels the statement \texttt{corrected\_verified} and says that notation, typography, and wording were normalized without changing the mathematical content. These are supplied editorial qualifications, not independent mathematical verification of this draft.
\end{itemize}
\end{document}
